\documentclass[a4paper,10pt]{report}\input{../0.Préambule/Préambule_cours_prof}\begin{document}

\pagestyle{empty}
\newpage
\noindent NOM, Prénom et classe : \dotfill

\section*{Devoir maison : Les fractions}
\addcontentsline{toc}{section}{Devoir maison : Les fractions}

\vspace{-1cm}
\begin{flushright}
Pour le vendredi 8 avril 2022
\end{flushright}

\paragraph{Exercice 1}: \hfill \textbf{10 points}

\medskip
L'Egypte est le premier producteur mondial de dattes. $\np{1700000}$ tonnes de dattes sont récoltées chaque année.

\medskip
\noindent Les quatre vingt-cinquièmes de cette production est exportée.
\begin{enumerate}
\item \begin{enumerate}
\item Quelle est la capitale de l'Égypte ?
\item Coller une image du drapeau égyptien.
\item Quel est l'emblème national de l'Égypte ?
\end{enumerate}
\item Combien de tonnes de dattes égyptiennes sont exportées chaque année ?
\item Combien de tonnes de dattes sont consommées par les Égyptiens ?
\item Il y a environ $102,3$ millions d'habitants en Égypte.

\noindent Combien chaque habitant consomme-t-il, en moyenne, de kg de dattes par an ?
\end{enumerate}

\medskip
\paragraph{Exercice 2}: \hfill \textbf{5 points}

\vspace{-8mm}
\noindent \begin{minipage}{11cm}

\vspace{3mm}
Écris la fraction de l'aire du carré que représente chaque morceau.

\medskip
\renewcommand{\arraystretch}{2}
\begin{tabularx}{\linewidth}{*{5}{>{\centering \arraybackslash}X}}
$A$ = \dotfill &
$B$ = \dotfill &
$C$ = \dotfill &
$D$ = \dotfill &
$E$ = \dotfill \\
\end{tabularx}

\end{minipage}
\begin{minipage}{5cm}
\begin{center}
\includegraphics[scale=0.3]{../40.Image/DM1 Les fractions.png}
\end{center}
\end{minipage}

\medskip
\paragraph{Exercice 3}: \hfill \textbf{5 points}

\medskip
\noindent \includegraphics[scale=0.54]{../40.Image/DM1 Les fractions 2.png}

\vspace{-12cm}
\noindent \begin{minipage}{8.5cm}
\medskip
\noindent Colorie d'une même couleur les cases égales.

\medskip
\noindent Quel est le nombre de la case non coloriée ?

\vspace{18mm}
\noindent \textbf{La suite est facultative ! Uniquement pour s'amuser  =) }


\medskip
\noindent Tu dois placer les dominos dans le parcours en les \\recopiant, sachant qu'un domino ne peut servir \\ qu'une seule fois.
\end{minipage}

\vspace{-2.5cm}
\noindent \Large{\ding{33}}  \dotfill
\end{document}